% Pursuing Stacks (cotes 134-2 to 134-8): archive cote 134-3, site page 27 (PDF page 28), typescript page 27.
% typescript. Draft compared: ps2.tex:1256-1314.
% AI-transcribed by Claude Opus 5.5 (claude-opus-5-5, adaptive thinking, effort high), Message Batches API, batch ps1 (msgbatch_01Rv5hCaNTtUt4JHx1yqoHPL), 2026-11-12; not yet checked by eye.
\scanpage{134-3/27}
\mspage{27}

Once this internal condition on B is understood,  step b) then would
amount to describing,  in terms of an arbitrary ``eligible'' algebraic
structure species \xout{xxxxx} expressed by the category M,  of what we should
mean by ``eligible functors'' \xout{BxxxMxxx} \hand{$(^{29})$}\margin{we will \unsure{rather say} ``test functors'', see below\dots}
\[ B \longrightarrow M \ , \]
giving rise in the usual way to a functor
\[ M \rightarrow B^{\wedge} \ . \]
In any case,  the latter functor defines upon M an induced \unsure{topology},
T say,  and the comparison \inserted{lemma} tells us that if either $B \rightarrow M$ \inserted{or $M \rightarrow B^{\wedge}$} is fully
faithful ,  then the topos associated to M is canonically equivalent\xout{x}
to $B^{\wedge}$  (using this comparison lemma for the functor which happens to be
fully faithful).  \xout{Onxthexotherxhandx} From this follows that the functor
(*)  $M^{\sim} \rightarrow$ (Hot)  defined by T \inserted{(cf. section 25)} is nothing but the compositum
\[ M^{\sim} \xrightarrow{\hand{$\approx$}} B^{\wedge} \longrightarrow (Hot) \quad , \]
and hence a localization functor.  In other words,  when B is a category
satisfying the condition seen above\hand{$(^{29})$}, then \uline{any} functor $B \rightarrow M$ satisfying\margin{a ``test category'', \unsure{as we will say}}
one of the two fullness conditions above \xout{xxx} yields a corresponding des-%
cription of (Hot) as a localization of $M^{\sim}$.  What is still lacking though
is a grasp on when two such functors  $B \rightarrow M$,  $B' \rightarrow M$  define essentially
``the same'' functor  $M \rightarrow$ (Hot),  or (more or less equivalently)  two
(Hot)-equivalent topologies T,T' on M;  \uline{is it enough,}  \uline{for instance,}
\uline{that they give rise to the same notion of} \uline{``weak equivalences''}  (namely\margin{\unsure{same question for Hot-equivalence of topologies on M}\dots} \hand{\illegible{1}}
morphisms in M which are tranformed into an isomorphism of (Hot)) ?
And moreover,  \xout{this} granting that this equivalence relation between
\inserted{certain} full subcategories \inserted{(say)} B of M is understood,  how define,in terms of M ,
a ``natural'' equivalence class of such full subcategories,  giving rise
to a \uline{canonical} functor  $M \rightarrow$ (Hot) \deleted{\hand{$(^{25})$}} \hand{?} \hand{$(^{25})$}

Recalling that the algebraic structure considered can be described
in terms of an arbitrary small category $\underline{B}$ where arbitrary finite direct
limits exist  (namely $\underline{B}$ is the full subcategory of M of objects of
finite presentation),  it seems reasonable to assume that indeed
\[ B \hookrightarrow \underline{B} \quad \text{(a full embedding)} \ , \]
and the question transforms into describing a\deleted{n} \inserted{natural} equivalence class \xout{xxxxx}
of such full subcategories,  in \inserted{(more or less)} any small category $\underline{B}$  where finite
direct limits exist,  and where moreover there exist such full subcate-%
gories B.  Also,  we may \deleted{to} have to throw in some extra\deleted{s} conditions on
$\underline{B}$ ,  such as the condition that direct sums be ``disjoint'' and ``universa\unsure{l}'' \hand{\unsure{$(^{*})$}}
already contemplated before.
